\gdef\mpxshipout{\shipout\hbox\bgroup
  \setbox0=\hbox\bgroup}
\gdef\stopmpxshipout{\egroup  \dimen0=\ht0 \advance\dimen0\dp0
  \dimen1=\ht0 \dimen2=\dp0
  \setbox0=\hbox\bgroup
    \box0
    \ifnum\dimen0>0 \vrule width1sp height\dimen1 depth\dimen2 
    \else \vrule width1sp height1sp depth0sp\relax
    \fi\egroup
  \ht0=0pt \dp0=0pt \box0 \egroup}
\mpxshipout% line 70 ./bimorph_fiziko.mp
fixed end%
\stopmpxshipout
\mpxshipout% line 77 ./bimorph_fiziko.mp
$\phi=100\,\rm V$%
\stopmpxshipout
\mpxshipout% line 78 ./bimorph_fiziko.mp
$\phi=0\,\rm V$%
\stopmpxshipout
\mpxshipout% line 85 ./bimorph_fiziko.mp
$E_y<0\;\;(\simeq-10^5\,\mathrm{V\,m^{-1}})$%
\stopmpxshipout
\mpxshipout% line 93 ./bimorph_fiziko.mp
$L=0.1\,\rm m$%
\stopmpxshipout
\mpxshipout% line 94 ./bimorph_fiziko.mp
$h_{\rm total}=1\,\rm mm$%
\stopmpxshipout
\mpxshipout% line 96 ./bimorph_fiziko.mp
drawn with height exaggerated $\times40$%
\stopmpxshipout
\mpxshipout% line 99 ./bimorph_fiziko.mp
top: $\epsilon^*_{xx}<0$ (contraction)%
\stopmpxshipout
\mpxshipout% line 101 ./bimorph_fiziko.mp
bottom: $\epsilon^*_{xx}>0$ (expansion)%
\stopmpxshipout
\mpxshipout% line 111 ./bimorph_fiziko.mp
orange: upper layer%
\stopmpxshipout
\mpxshipout% line 114 ./bimorph_fiziko.mp
teal: lower layer%
\stopmpxshipout
\mpxshipout% line 116 ./bimorph_fiziko.mp
schematic piezoelectric bimorph%
\stopmpxshipout
\end{document}
